\documentclass[11pt]{article}
\usepackage{amsmath,amssymb,amsthm}
\usepackage[margin=1in]{geometry}
\usepackage{hyperref}
\newtheorem{theorem}{Theorem}
\newtheorem{lemma}{Lemma}
\newcommand{\Q}{\mathbb Q}
\newcommand{\Z}{\mathbb Z}
\newcommand{\tors}{\mathrm{tors}}
\newcommand{\Jac}{\operatorname{Jac}}
\newcommand{\im}{\operatorname{im}}
\title{Infinitely many split genus-2 Jacobians with rational torsion
  \(\Z/4\Z\times\Z/16\Z\)}
\author{}
\date{20 October 2026}
\begin{document}
\maketitle

\begin{abstract}
We give an explicit rational family of smooth genus-2 curves over
\(\Q\) whose Jacobians split over \(\Q\) and have rational torsion
subgroup exactly \(\Z/4\Z\times\Z/16\Z\).
The family contains infinitely many geometric isomorphism classes.
Exactness follows from the intersection of the torsion Kummer images
of the two elliptic factors.
\end{abstract}

\section{The family}

For \(s\in\Q\), put
\begin{equation}\label{eq:t}
t=t(s)=-
\frac{3(3s^2-10s+5)^2+(3s^2-5)^2}{4(3s^2-5)^2},
\qquad D(t)=256t^2+47t+16.
\end{equation}
Let \(C_s\) be the smooth projective model of
\begin{equation}\label{eq:family}
y^2=D(t)(16x^2+9t)
       \bigl(175x^2-9(4t-1)^2\bigr)
       \bigl(9x^2+(4t+1)^2\bigr).
\end{equation}

\begin{theorem}\label{thm:main}
For every \(s\in\Q\), the curve \(C_s\) has genus \(2\), its Jacobian
is isogenous over \(\Q\) to a product of elliptic curves, and
\[
 \Jac(C_s)(\Q)_{\tors}\simeq\Z/4\Z\times\Z/16\Z.
\]
The curves \(C_s\), with \(s\in\Q\), contain infinitely many
\(\overline{\Q}\)-isomorphism classes.
\end{theorem}

We prove exactness for every rational parameter. In particular,
the theorem requires no exceptional thin set or congruence restriction.

\section{Parameters and elliptic gluing}

Put
\[
 a=\frac{3(3s^2-10s+5)}{2(3s^2-5)},\qquad
 b=-\frac{5(3s^2-6s+5)}{2(3s^2-5)}.
\]
A direct calculation gives
\begin{equation}\label{eq:conic}
 a^2=-3(t+1/4),\qquad b^2=-5(t-1/4),\qquad
 6b^2-10a^2=15.
\end{equation}
The polynomials \(3s^2-5\), \(3s^2-10s+5\), and \(3s^2-6s+5\)
have no rational zeros. Hence \(a,b\ne0\) and \(t<-1/4\) for every
rational \(s\). The discriminant of \(D\) is \(-14175\), so
\(D(t)>0\) for real \(t\).

Let \(m=t^2-1/16\), and consider
\begin{align*}
 E_t &: y^2+xy-my=x^3-mx^2,\\
 F &: y^2+\frac{47}{42}xy+\frac{5}{49}y
                 =x^3+\frac{5}{49}x^2.
\end{align*}
After completing the square in \(y\), their monic cubic models have
the following ordered roots:
\begin{equation}\label{eq:roots}
 \begin{split}
 (e_1,e_2,e_3)&=(m,-t/2-1/8,t/2-1/8),\\
 (f_1,f_2,f_3)&=(-15/112,-6/49,-10/63).
 \end{split}
\end{equation}
The roots are distinct for \(t<-1/4\).
Let \(\alpha:E_t[2]\longrightarrow F[2]\) send the point at \(e_i\)
to the point at \(f_i\), and let \(\Gamma\) be its graph.

We apply the gluing construction of Howe--Lepr\'evost--Poonen
\cite[Propositions 3 and 4]{HLP}. In their notation,
\[
 a_2=\sum_{i=1}^3 e_i(f_{i+2}-f_{i+1})=-\frac{D(t)}{7056}\ne0,
\]
where subscripts are taken modulo \(3\). Thus the correspondence of
roots is not induced by an affine change of the elliptic \(x\)-coordinate,
and \(\alpha\) is not the restriction of an isomorphism of elliptic curves
over \(\overline{\Q}\).
Consequently
\[
 J=(E_t\times F)/\Gamma
\]
with its induced principal polarization is a genus-2 Jacobian.
Write
\[
 \phi:E_t\times F\longrightarrow J,\qquad
 \psi:J\longrightarrow E_t\times F
\]
for the gluing isogeny and its dual, so that
\(\psi\phi=[2]\) and \(\phi\psi=[2]\).

The quotient is precisely the Jacobian of \eqref{eq:family}.
For completeness, let \(h(X)\) be the sextic in
\cite[Proposition 4]{HLP}, using the root order \eqref{eq:roots}.
Its constants \(A,B\) are
\[
 A=-\frac{25D(t)}{17076987648},\qquad
 B=\frac{t(4t-1)^2(4t+1)^2D(t)}{1806336}.
\]
If \(g_t(x)\) denotes the right-hand side of \eqref{eq:family}, then
\begin{equation}\label{eq:scale}
 h(7x)=\left(\frac{Atm^2}{12096}\right)^2g_t(x).
\end{equation}
All factors in this change of variables are nonzero for our parameters.
The sextic discriminant is
\begin{equation}\label{eq:disc}
 \operatorname{disc}_x(g_t)=
 56171432824012800\,t(4t-1)^2(4t+1)^2D(t)^{22}\ne0.
\end{equation}
Thus every \(C_s\) is smooth of genus \(2\), and its Jacobian splits
over \(\Q\).

\section{The torsion Kummer images}

We use multiplicative notation for \(\Q^\times/\Q^{\times2}\).
On an elliptic curve with monic cubic model
\(Y^2=\prod_{i=1}^3(x-r_i)\), the \(2\)-Kummer map is
\[
 \delta:E(\Q)/2E(\Q)\longrightarrow
          (\Q^\times/\Q^{\times2})^3.
\]
For a point away from the nonzero \(2\)-torsion, it is represented by
\((x-r_1,x-r_2,x-r_3)\).
At the \(2\)-torsion point with \(x=r_i\), the \(i\)-th component is
\(\prod_{j\ne i}(r_i-r_j)\), and the other components are \(r_i-r_j\).
The identity has image \((1,1,1)\).
We use \eqref{eq:roots} for the coordinate order on the two curves,
so that their Kummer coordinates agree under \(\alpha\).

Let \(P=(0,0)\) on the displayed generalized Weierstrass model of \(E_t\),
and let \(Q=(0,0)\) on that of \(F\).
Let \(T\in E_t[2]\) have \(x=e_2\), and let \(U\in F[2]\)
have \(x=f_3\).
The addition law gives
\[
 \operatorname{ord}(P)=4,\quad x(2P)=e_1,\qquad
 \operatorname{ord}(Q)=8,\quad x(4Q)=f_2.
\]
In particular, \(P,T\) generate a subgroup of type \([2,4]\),
and \(Q,U\) generate a subgroup of type \([2,8]\).

The Kummer images, as triples of square classes, are
\begin{align}
 \delta_E(P)&=(-m,\,2(t+1/4),\,-2(t-1/4)),&
 \delta_E(T)&=(-1,t,-t),\label{eq:E-kummer}\\
 \delta_F(Q)&=(105,6,70),&
 \delta_F(U)&=(-7,-1,7).\label{eq:F-kummer}
\end{align}
By \eqref{eq:conic}, \(m=a^2b^2/15\). Therefore
\begin{equation}\label{eq:match}
 \delta_E(P)=(-15,-6,10)=\delta_F(Q+U).
\end{equation}

\begin{lemma}\label{lem:parents}
For every rational \(s\),
\[
 E_t(\Q)_{\tors}=\langle P,T\rangle\simeq
        \Z/4\Z\times\Z/2\Z,\qquad
 F(\Q)_{\tors}=\langle Q,U\rangle\simeq
        \Z/8\Z\times\Z/2\Z.
\]
\end{lemma}
\begin{proof}
Mazur's theorem \cite{Mazur} gives the assertion for \(F\).
For \(E_t\), the only possibilities are \([2,4]\) and \([2,8]\).
In the latter case its four points of order \(4\) are
\(\pm P,\pm(P+T)\), and some such point is twice a rational point of
order \(8\).
However, the first Kummer coordinates of \(P\) and \(P+T\)
have square classes \(-15\) and \(15\), respectively.
Neither is a square in \(\Q\). Their negatives have the same Kummer
images, so none of these points is divisible by \(2\) in \(E_t(\Q)\).
This excludes \([2,8]\).
\end{proof}

Put \(T_E=E_t(\Q)_{\tors}\) and \(T_F=F(\Q)_{\tors}\).
The first coordinates of their Kummer images are, respectively,
\[
 \{1,-15,-1,15\},\qquad \{1,105,-7,-15\}.
\]
The two sets have exactly the common square classes \(1\) and \(-15\).
By \eqref{eq:match}, the complete triples with first coordinate
\(-15\) also agree. It follows that
\begin{equation}\label{eq:intersection}
 \delta_E(T_E)\cap\delta_F(T_F)
       =\{(1,1,1),(-15,-6,10)\}.
\end{equation}

\section{Exact torsion of the Jacobian}

Define
\[
 M=\{(A,B)\in T_E\times T_F:\delta_E(A)=\delta_F(B)\}.
\]
We claim that there is an exact sequence
\begin{equation}\label{eq:exact}
 0\longrightarrow\ker(\psi)(\Q)
 \longrightarrow J(\Q)_{\tors}
 \xrightarrow{\ \psi\ }M\longrightarrow0.
\end{equation}
Indeed, choose halves \(\widetilde A,\widetilde B\) of \(A,B\).
The point \(\phi(\widetilde A,\widetilde B)\) is rational if and only if
every Galois difference of the pair \((\widetilde A,\widetilde B)\)
lies in \(\Gamma\).
Since \(E_t[2]\) and \(F[2]\) are rational, this is exactly the equality
of Kummer classes in the definition of \(M\).
It is then a lift of \((A,B)\) under \(\psi\).
The lift is torsion: either this follows from its construction using
halves of torsion points, or from \(2R=\phi\psi(R)\).
Conversely, every rational lift gives the equality of Kummer classes.
This proves \eqref{eq:exact}.

The group scheme
\[
 \ker(\psi)=\phi(E_t[2]\times F[2])
       \simeq(E_t[2]\times F[2])/\Gamma
\]
is constant of order \(4\).
The kernel of the Kummer map on the torsion subgroup of an elliptic
curve is twice that torsion subgroup, since a half of a torsion point
is torsion.
Hence Lemma~\ref{lem:parents} and \eqref{eq:intersection} give
\[
 \#M=\#(2T_E)\#(2T_F)
       \#\bigl(\delta_E(T_E)\cap\delta_F(T_F)\bigr)
      =2\cdot4\cdot2=16.
\]
By \eqref{eq:exact},
\begin{equation}\label{eq:order}
 \#J(\Q)_{\tors}=4\cdot16=64.
\end{equation}

It remains to determine the group structure. Put \(V=Q+U\), which
has order \(8\), and
\[
 X=\phi(P,0),\qquad Y=\phi(0,V).
\]
The graph \(\Gamma\) is generated by \((2P,U+4V)\) and \((T,4V)\).
In the quotient of \(T_E\times T_F\), these relations eliminate
the generators corresponding to \(T\) and \(U\).
The remaining relations are \(4X=0\) and \(8Y=0\).
Thus
\[
 H:=\phi(T_E\times T_F)=\langle X,Y\rangle
       \simeq\Z/4\Z\times\Z/8\Z.
\]
The match \eqref{eq:match} gives \(R\in J(\Q)_{\tors}\)
with \(\psi(R)=(P,V)\). It satisfies
\[
 2R=X+Y.
\]
The point \(X+Y\) has order \(8\), so \(R\) has order \(16\).
Moreover, \(R\notin H\), since \(2H\) has exponent \(4\).
Therefore \(\langle H,R\rangle\) has order \(64\), and by
\eqref{eq:order} it is the whole rational torsion subgroup.
It is generated by \(X,R\), since \(Y=2R-X\).
The resulting surjection
\[
 \Z/4\Z\times\Z/16\Z\longrightarrow J(\Q)_{\tors}
\]
is an isomorphism because both groups have order \(64\).

\section{Variation in moduli}

The varying elliptic factor has
\[
 j(E_t)=
 \frac{(256t^4+224t^2+1)^3}
      {t^2(4t-1)^4(4t+1)^4}.
\]
This rational function is nonconstant, and so is \(t(s)\).
Thus \(j(E_{t(s)})\) assumes infinitely many values for \(s\in\Q\).
The gluing construction supplies a degree-\(2\) map
\(C_s\longrightarrow E_t\) over \(\Q\).
A fixed smooth genus-2 curve over \(\overline{\Q}\) has finitely
many automorphisms, and hence finitely many involutions and
degree-\(2\) elliptic quotients up to isomorphism.
Consequently the infinitely many values of \(j(E_{t(s)})\)
force infinitely many geometric isomorphism classes of \(C_s\).
This proves Theorem~\ref{thm:main}.

\section*{Reproducibility and the census example}

The accompanying standalone Magma file \texttt{family\_416\_split\_certificate.m} verifies
the elliptic point orders, all Kummer identities, the HLP change of
variables \eqref{eq:scale}, the discriminant \eqref{eq:disc},
the conic parametrization, and the elliptic \(j\)-formula.
All assertions passed in Magma V2.30-1.
The proof above uses only exact identities, elliptic torsion
classification, and the stated gluing theorem.

At \(s=1\), we have \(t=-1\), and after dividing \(y\) by \(225\)
the curve becomes
\[
 y^2=(16x^2-9)(7x^2-9)(x^2+1).
\]
The change of variables \(z=3/x\), \(w=3y/x^3\) gives
\[
 w^2=z^6-14z^4-95z^2+1008.
\]
Writing \(w=2v+z^3+z\) gives the census example \texttt{4.16.a},
\[
 v^2+(z^3+z)v=-4z^4-24z^2+252.
\]
Magma independently computes exact torsion \([4,16]\) for this
specialization and for \(s=2\), where \(t=-19/49\).
These specialization checks are consistency checks; they are not
needed for the uniform torsion argument.

\begin{thebibliography}{9}
\bibitem{HLP}
E.~W. Howe, F. Lepr\'evost, and B. Poonen,
\emph{Large torsion subgroups of split Jacobians of curves of genus two or three},
Forum Math. \textbf{12} (2000), 315--364.
\url{https://math.mit.edu/~poonen/papers/large.pdf}.

\bibitem{Mazur}
B. Mazur,
\emph{Modular curves and the Eisenstein ideal},
Publ. Math. IHES \textbf{47} (1977), 33--186.
\url{https://www.numdam.org/item/PMIHES_1977__47__33_0/}.
\end{thebibliography}
\end{document}
